This chapter covers quad assembly and mechanical integration of the wearable enclosure.

%------------------------------------------------------------------------------
\subsection{Drone Mechanical Design}
\label{subsec:drone}
%------------------------------------------------------------------------------

Two ground-station suites (Betaflight Configurator and Mission Planner, depending on firmware stage)
were used to configure motor spin, calibrate ESC endpoints, verify receiver channels, and complete
pre-flight checklists.

Motor direction and propeller pairing received particular attention: each corner was tested in
isolation because incorrect torque or blade orientation leads to immediate flips on takeoff.
After firmware setup, motors were spun in the Configurator and visually checked before props were
mounted.
Clockwise and counter-clockwise props were installed according to the Mark4-style frame geometry,
then throttle sweeps confirmed all four corners produced thrust in the correct direction.

Those steps reduced crash risk during the intensive flight campaign and supported repeated takeoff,
hover, and translation tests invoked from either the Pocket transmitter or the wristband.

Beyond rigging, the airframe is a five-inch class FPV build on a Mark4-style kit with four 1860~KV
brushless motors, five-inch props, and a four-piece prop guard set for low-altitude experiments.

%------------------------------------------------------------------------------
\subsection{Wristband Enclosure Design}

\subsubsection{Enclosure Requirements and CAD Design}

A dedicated wrist-wearable enclosure was CADed for bench and field demos.
Internal bosses position the ESP32, VC-02 board, ELRS module, and wiring harness; cutouts expose the
OLED window and the microphone port so keywords remain intelligible while the display stays visible.
The form factor targets low mass and short moment arm on the pilot's wrist.

\subsubsection{3D Printing and Assembly}

The housing was manufactured with consumer FDM printers across several fit-check iterations.
After tuning tolerances, the stack became more compact than the early breadboard prototype and
survived repeated strap-on operations during outdoor flights.
